\documentclass[10pt,a4paper]{amsart}
\usepackage[margin=19mm]{geometry}
\usepackage{amsmath,amssymb,mathtools}
\usepackage[scr=rsfs,cal=euler]{mathalfa}
\usepackage{xfrac}
\usepackage[unicode,hidelinks,bookmarks=false]{hyperref}
\newcommand{\ie}{i.e.\ }
\newcommand{\cf}{cf.\ }
\newcommand{\resp}{respectively\ }
\newcommand{\Subsection}{\subsection}
\providecommand{\nobreakdash}{\nobreak\hbox{-}}
\setlength{\emergencystretch}{2em}
\multlinegap0pt
\usepackage{booktabs}
\usepackage{multirow}
\usepackage{bm}
\usepackage{algorithm}\usepackage{algpseudocode}
\usepackage{caption}
\usepackage{xcolor}
\newtheorem{theorem}{Theorem}
\newcommand{\T}{^\top}
\newcommand{\tvd}{d_\textrm{TV}}
\title{Safety Beyond the Training Data: Robust Out-of-Distribution MPC via Conformalized System Level Synthesis}
\author{Anutam Srinivasan, Antoine Leeman, Glen Chou}
\date{Source: arXiv:2602.12047v1 (CC BY 4.0)}
\begin{document}
\maketitle

\noindent\emph{Source and licence.} Safety Beyond the Training Data: Robust Out-of-Distribution MPC via Conformalized System Level Synthesis. Version arXiv:2602.12047v1 (CC BY 4.0), distributed by arXiv under the Creative Commons Attribution 4.0 International licence, http://creativecommons.org/licenses/by/4.0/, as recorded in the arXiv licence metadata for this exact version and shown on the abstract page https://arxiv.org/abs/2602.12047v1. Full licence notice: \texttt{licenses/arxiv-2602.12047-CC-BY-4.0.txt}.\par\smallskip

\noindent\textit{Editorial scope.} Counted unit: the article's theorem thm:tube\_gap (source0.tex lines 1153--1159). The article's complete original proof of it is delivered with it (source0.tex lines 1426--1444, one \texttt{proof} environment of 19 lines, which the article places away from the statement). No mathematics is written, repaired or re-derived: the statement and the proof are the article's own lines, in the article's own notation, and no retained line is substituted or re-flowed.  Retained uncounted as the source-local context the proof needs: lines 1129--1131; lines 1138--1145; lines 1377--1379; lines 1403--1408.  Nothing was omitted from any retained span, so the package carries no omission record.  No retained line contains a citation, so the wrapper carries no bibliography and no bibliography entry.  No retained line contains a figure environment, an image-inclusion command or drawing code, so no image is delivered and the package declares no asset dependency.  Editorial formatting only, in addition: an AMS \texttt{amsart} preamble that restates the article's own declarations and macros used by the retained lines, namely the environments \texttt{theorem} on separate counters, together with the standard packages those lines need.  The retained environments print in this document as Theorem 1 in order of appearance, whereas the article prints the same items with the numbering of its own full document.  The complete native source of the article as distributed in the arXiv e-print for this exact version is bundled unchanged as \texttt{sources/194517/source0.tex}.
\begin{equation}\label{eq:lipschitz}
\tvd( S_{i,k},S_{k,k}) \leq  \epsilon ||[z_k\T,v_k\T] - [x_i\T,u_i\T]||_{2},    
\end{equation}

\begin{theorem}\label{thm:cov_gap_thm} The \textit{Coverage Gap}$\,=\sum_{i=1}^{N_\textrm{calib}}\tilde{w}_i\,\tvd(S^{i,k}, S^{k,k})$ of our conformal predictor satisfies

\vspace{-10pt}
\begin{equation}\textstyle
    \text{Coverage Gap}\le \underbrace{\textstyle\sum_{i=1}^{N_\textrm{calib}} \left( \frac{\rho^{d_i}}{1 + \sum_{j=1}^{N_\textrm{calib}} \rho^{d_j}} \right) \cdot 2\epsilon \cdot d_i}_{\text{Tight Bound}} \le \underbrace{\textstyle 2 \epsilon \left[  \frac{d_1 }{1 - \rho^{d_\text{min}}} + \frac{d_\text{max}\rho^{d_\text{min}}}{(1-\rho^{d_\text{min}})^2}\right]}_{\text{Interpretable Bound}}, 
\label{eq:thm3_bound}\end{equation}
where we sort the distance to points in the calibration set such that $d_i = ||[z_k\T,v_k\T] - [x_i\T,u_i\T]||_{2}$, $d_1 \le d_2 \le \dots \le d_n$, $d_\text{min} = \min_{1 \leq i \leq n-1}\{d_{i+1} - d_i\}$, and $d_\text{max} = \max_{1 \leq i \leq n-1} \{ d_{i+1} - d_i\}$.
\end{theorem}

\begin{equation}\label{eq:init_barber}
\text{Coverage Gap} \le \sum_{i=1}^{N_\textrm{calib}} \tilde{w}_i \cdot d_{TV}(S^{i,k}, S^{k,k}) \le \sum_{i=1}^{N_\textrm{calib}} \tilde{w}_i \cdot 2 d_{TV}(S_{i,k}, S_{k,k})
\end{equation}

\begin{align}
    &\text{Coverage Gap} \le 2\epsilon\frac{N_\text{sum}}{D_\text{sum}} \\
    &\le 2 \epsilon \left[{d_1 \left( \frac{1 - \rho^{N_\textrm{calib} \cdot d_\text{min}}}{1 - \rho^{d_\text{min}}} \right) +
    d_\text{max} \left( \frac{\rho^{d_\text{min}} \overbrace{- N_\textrm{calib} \rho^{N_\textrm{calib} \cdot d_\text{min}} + (N_\textrm{calib}-1)\rho^{(N_\textrm{calib}+1)d_\text{min}}}^{\leq 0}}{(1-\rho^{d_\text{min}})^2} \right)}\right]\underbrace{\frac{1 - \rho^{d_\text{max}}}{1 - \rho^{N_\textrm{calib} \cdot d_\text{max}}}}_{\leq 1}\label{eq:n_sum_d_sum}\\
    &\le 2 \epsilon \left[  \frac{d_1 }{1 - \rho^{d_\text{min}}} + \frac{d_\text{max}\rho^{d_\text{min}}}{(1-\rho^{d_\text{min}})^2}\right].\label{eq:final_bound} 
\end{align}

\begin{theorem}\label{thm:tube_gap}
Under the same assumption of Theorem \ref{thm:cov_gap_thm}, we have that for all $(x,u) \in \mathcal{R}_k$, the tube around the nominal point $(z_k,v_k)$, 
% \begin{equation}\textstyle
    $\Pr[\varepsilon(x,u) \in \mathcal{E}(z_k,v_k)] \geq 1-\alpha_k- 2\sum_{i=1}^{n} \tilde{w}_i d_{TV}(S_{i,k}, S_{k,k})- \gamma({\mathcal{R}_k})$,
% \end{equation}
for $ \gamma({\mathcal{R}_k}) := 2\hat{\epsilon}M(\mathcal{R}_k) $, where $M(\mathcal{R}_k)$ is the maximum axis length of tube $\mathcal{R}_k$, and $\hat \epsilon \le \epsilon$ is chosen such that \eqref{eq:lipschitz} holds for all $(x, u) \in \mathcal{R}_k$.%, and is less than the global Lipschitz constant in $\epsilon$, Theorem \ref{thm:cov_gap_thm}.
\end{theorem}

\begin{proof}\textbf{(Theorem \ref{thm:tube_gap})}
    Using Equation \ref{eq:init_barber}, we have that 
    \begin{equation}\label{eq:init_barber_2}
    \text{Coverage Gap} \le \sum_{i=1}^{N_\textrm{calib}} \tilde{w}_i \cdot d_{TV}(\tilde{S}^{i,k}, \tilde{S}^{k,k}) \le \sum_{i=1}^{N_\textrm{calib}} \tilde{w}_i \cdot 2 d_{TV}(S_{i,k}, \tilde{S}_{k,k}),
    \end{equation}
    where $(s_{1,k}, s_{2,k}, s_{i-1,k}, \tilde{s}_{k,k}, s_{i+1,k}, \ldots, s_{N_\textrm{calib},k}) \sim \tilde{S}^{i,k}$, $(s_{1,k}, \ldots, s_{N_\textrm{calib},k}, \tilde{s}_{k,k}) \sim \tilde{S}^{k,k}$, $\tilde{s}_{k,k} = s_k(x_k,u_k,f(x_k,u_k))$ is unknown due to uncertainty in $f(x_k,u_k)$, and $\tilde{s}_{k,k} \sim \tilde{S}_{k,k}$ is the distribution of the non-conformity score corresponding to the observed state and control $(x_k,u_k)$. Then, using the triangle inequality we have that,
    \begin{equation}\label{eq:split_into_sums}
    \text{Coverage Gap} \le \sum_{i=1}^{N_\textrm{calib}} \tilde{w}_i \cdot d_{TV}(\tilde{S}^{i,k}, \tilde{S}^{k,k}) \le \underbrace{\sum_{i=1}^{N_\textrm{calib}} \tilde{w}_i \cdot 2 d_{TV}(S_{i,k}, S_{k,k})}_{(*), \text{Bounded in Eq. \ref{eq:final_bound}}} + \underbrace{\sum_{i=1}^{N_\textrm{calib}} \tilde{w}_i \cdot 2 d_{TV}(S_{k,k}, \tilde{S}_{k,k})}_{(**)}.
    \end{equation}
    Observing that the left sum $(*)$ is already bounded by Theorem \ref{thm:cov_gap_thm}, we proceed to bound the right sum $(**)$. Then, realizing the sum of normalized weights is $\le 1$, we obtain, 
    \begin{equation}\label{eq:near_by_cov}
        \sum_{i=1}^{N_\textrm{calib}} \tilde{w}_i \cdot 2 d_{TV}(S_{k,k}, \tilde{S}_{k,k}) = 2 d_{TV}(S_{k,k}, \tilde{S}_{k,k})\sum_{i=1}^{N_\textrm{calib}} \tilde{w}_i \le 2 d_{TV}(S_{k,k}, \tilde{S}_{k,k})\
    \end{equation}
     Then, using the Lipschitz-type bound, 
    \begin{equation}
        2 d_{TV}(S_{k,k}, \tilde{S}_{k,k}) \le 2\hat{\epsilon}||[x_k\T,u_k\T]-[z_k\T,v_k\T]||_2 \le 2\hat{\epsilon}M(\mathcal{R}_k) =: \gamma({\mathcal{R}_k}),
    \end{equation} 
    where, $(x_k,u_k) \in \mathcal{R}_k$, $M(\mathcal{R}_k)$ is the maximum axis length of the tube, and $\hat{\epsilon} \le \epsilon$ is the local Lipschitz constant in the tube. The remaining proof follows substituting the bound, $\gamma({\mathcal{R}}_k)$ for $(**)$ in \eqref{eq:split_into_sums}.
\end{proof}

\end{document}
